% Background. Model definitions prepare the contributions and their interpretation.
\section{Background}
% Source. Background 2.1, eq:class-probability-vector, eq:categorical-likelihood,
% eq:likelihood-factorization, eq:bayes-posterior, eq:bma-predictive.
% Speaking notes
% We consider C classes indexed by c, input x, and a generic label y.
% The training data contain N input and label pairs, indexed by i. The observed label for example i is y_i.
% No component subscript denotes a whole vector. Subscript c selects class c. Subscript y selects the component for label y.
% For given weights w, the network deterministically outputs the vector f of class probabilities.
% The displayed conditions say that each probability is nonnegative and the probabilities sum to one.
% A vector with these properties lies on the probability simplex. We state the conditions directly on the slide.
% The categorical distribution assigns each possible label its corresponding component of f.
% For training example i, we evaluate it at the observed label y_i. This selects f_{y_i}(x_i;w).
% For example, if y_i is 3, this selects f_3(x_i;w), the probability assigned to class 3.
% We assume the observations are independent given their inputs and the shared weights.
% Multiplying their categorical probabilities gives the likelihood of all observed training labels.
% The prior describes our assumptions about the weights before observing these labels. This is normally taken to be isotropic Gaussian. We also use isotropic gaussian in our experiments.
% Combining prior and likelihood gives the posterior over weights, up to a normalising constant.
% We retain uncertainty about weights and average predictions over this posterior.
% Next, we consider a network that outputs a distribution over class probabilities.
\begin{frame}{Categorical Bayesian Neural Networks}
\label{frm:bg-categorical-bnn}
\setlength{\abovedisplayskip}{4pt}
\setlength{\belowdisplayskip}{4pt}
\setlength{\abovedisplayshortskip}{4pt}
\setlength{\belowdisplayshortskip}{4pt}
\begin{itemize}
\setlength{\itemsep}{4mm}
\item $C$ classes, input $x$, label $y\in\{1,\ldots,C\}$,
training data $D=\{(x_i,y_i)\}_{i=1}^{N}$.

\item The network with weights $w$ outputs class-probability vector,
% Thesis eq:class-probability-vector and simplex constraints.
\[
f(x;w)=\bigl(f_1(x;w),\ldots,f_C(x;w)\bigr),
\quad f_c(x;w)\geq0,\quad \sum_{c=1}^{C}f_c(x;w)=1.
\]

\item Categorical likelihood and weight prior define the posterior over weights.
% Thesis eq:likelihood-factorization with eq:categorical-likelihood substituted,
% and eq:bayes-posterior without its normalising constant.
\begin{align*}
p(D\mid w)&=\prod_{i=1}^{N}\mathrm{Cat}\bigl(y_i\mid f(x_i;w)\bigr)
=\prod_{i=1}^{N}f_{y_i}(x_i;w),\\[2mm]
\underbrace{p(w\mid D)}_{\text{posterior}}
&\propto\underbrace{p(D\mid w)}_{\text{likelihood}}\,
\underbrace{p(w)}_{\text{prior}}.
\end{align*}
\end{itemize}
\end{frame}

% Source. Background 2.2, eq:dirichlet-density, eq:general-dirichlet-mean,
% eq:dirichlet-variance. The probability vectors represent class probabilities in our classification setting.
% Speaking notes
%“In the categorical BNN, a given input and network weights produce one class-probability vector \(f(x;w)\). We now introduce a distribution over possible class-probability vectors, denoted \(\pi\). We use the Dirichlet distribution for this.”
% In our classification setting, pi is a random class-probability vector. Its component pi_c is the probability of class c.
% Each possible vector has nonnegative entries that sum to one.
% We model its distribution with a Dirichlet distribution, specified by positive concentration parameters alpha.
% The mean vector is m, with components m_c. Each m_c is the average of pi_c under Dir(pi|alpha).
% Its mean is alpha divided by the sum of its entries. We call that sum the total concentration parameter.
% The variance describes how each component of pi varies around its mean.
% At a fixed mean, larger total concentration gives smaller variance. this will be important later when we discuss the allocation of uncertainty.
% Next, we connect this distribution to the network output and the categorical label model.
\begin{frame}{The Dirichlet Distribution}
\label{frm:bg-dirichlet-distribution}
\setlength{\abovedisplayskip}{4pt}
\setlength{\belowdisplayskip}{4pt}
\begin{itemize}
\setlength{\itemsep}{4mm}
\item Random class-probability vector $\pi=(\pi_1,\ldots,\pi_C)$,
$\pi_c\geq0$ $\forall c$, $\sum_{c=1}^{C}\pi_c=1$.

\item $\pi$ is Dirichlet distributed with concentration parameters $\alpha=(\alpha_1,\ldots,\alpha_C)$, $\alpha_c>0$ $\forall c$,
% Thesis eq:dirichlet-density and its parameter domain.
\[
p(\pi\mid\alpha)=\mathrm{Dir}(\pi\mid\alpha).
\]

\item Mean vector and random class-probability variance,
% Thesis eq:general-dirichlet-mean and eq:dirichlet-variance.
\[
\begin{aligned}
&m=\mathbb E[\pi]=\frac{\alpha}{\alpha_0},
\qquad \underbrace{\alpha_0=\sum_{c=1}^{C}\alpha_c,}_{\text{total concentration parameter}}\\[4mm]
&\operatorname{Var}(\pi_c)=\frac{m_c(1-m_c)}{\alpha_0+1}.
\end{aligned}
\]
\begin{itemize}
\item At a fixed mean, larger total concentration gives less spread.
\end{itemize}
\end{itemize}
\end{frame}

% Source. Background 2.4 and 2.9, eq:dirichlet-classification-label,
% eq:fixed-weight-dirichlet and eq:fixed-weight-dirichlet-marginal.
% Methods 4.2.1, eq:edl-softplus-evidence and eq:edl-dirichlet-concentration.
% Speaking notes
% We want label probabilities for an input x.
% The network outputs alpha(x;w), specifying a Dirichlet distribution over class-probability vectors pi.
% Given pi, we model the label with a categorical distribution, so p(y|pi) equals pi_y.
% Together, these distributions define the Dirichlet-categorical model.
% The vector pi is unobserved in this model. We marginalise it to obtain p(y|x,w).
% The integral runs over all probability vectors in the simplex.
% Substituting p(y|pi) = pi_y gives the expectation of pi_y under the Dirichlet distribution.
% This is the corresponding mean component, alpha_y(x;w) divided by alpha_0(x;w).
% Thus the Dirichlet mean gives the label probabilities for x. (For one predicted class, choose their largest component.)
% We adopt the evidence parameterisation of evidential deep learning, or EDL, motivated by subjective logic.
% Subjective logic represents support for each class and a remaining mass not assigned to any class.
% EDL learns nonnegative evidence e_c for each class and maps it to Dirichlet parameters through alpha_c = e_c + 1.
% The +1 supplies one baseline unit per class. Zero evidence therefore gives the uniform Dirichlet Dir(1,...,1).
% This baseline is EDL's choice within subjective logic. A general Dirichlet distribution only requires alpha_c > 0.
% Our softplus output gives e_c > 0, so alpha_c > 1 and zero evidence is a limiting case.
% We adopt this parameterisation and later add a posterior over network weights. Evidence is a learned output, not an observed count.
% Source. Sensoy et al. 2018, Section 3, Eq. (1) and the zero-evidence example on page 4.
% We can now explain how the model decomposes predictive uncertainty.
\begin{frame}{Dirichlet Networks for Classification}
\label{frm:bg-dirichlet-network}
\setlength{\abovedisplayskip}{6pt}
\setlength{\belowdisplayskip}{6pt}
\begin{itemize}
\setlength{\itemsep}{7mm}
\item Dirichlet-categorical network outputs $\alpha(x;w)$. 
% Thesis eq:fixed-weight-dirichlet and eq:dirichlet-classification-label.
\begin{align*}
p(\pi\mid x,w)&=\mathrm{Dir}\bigl(\pi\mid\alpha(x;w)\bigr),\\
p(y\mid\pi)&=\mathrm{Cat}(y\mid\pi)=\pi_y.
\end{align*}
\item Marginalise $\pi$ to obtain label probabilities for input $x$.
% Thesis eq:dirichlet-classification-marginal and eq:fixed-weight-dirichlet-marginal.
\[
\begin{aligned}
p(y\mid x,w)
&=\int p(y\mid\pi)\,p(\pi\mid x,w)\,d\pi\\
&=\mathbb E_{\mathrm{Dir}(\pi\mid\alpha(x;w))}[\pi_y]
=\frac{\alpha_y(x;w)}{\alpha_0(x;w)}
=m_y(x;w).
\end{aligned}
\]
\item \textbf{Evidential deep learning (EDL)} parameterises the concentrations through evidence $e_c(x;w)>0$, with %we use softplus for our evidence, sensoy proposed relu. But we used edl pytorch library, which used softplus, maybe because its more stable?
% Methods eq:edl-softplus-evidence and eq:edl-dirichlet-concentration.
$\alpha_c(x;w)=e_c(x;w)+1$. % 
\end{itemize}
\end{frame}

% Source. Background 2.6 to 2.9, eq:expected-conditional-entropy,
% eq:mutual-information, eq:three-way-decomposition,
% eq:aleatoric-epistemic-decomposition and eq:aleatoric-distributional-split.
% Speaking notes
% Predictive uncertainty is the entropy of the model's predicted label distribution for an input.
% Aleatoric uncertainty is the expected conditional entropy obtained by conditioning on the class-probability vector.
% It measures the label uncertainty remaining once that vector is known.
% The vector is f(x;w) for the categorical BNN and pi for Dirichlet models.
% In the categorical BNN, average the categorical label entropy over the weight posterior.
% In the deterministic Dirichlet network, average it over the Dirichlet distribution of pi.
% Epistemic uncertainty is mutual information between the label y and weights w.
% It measures the reduction in label uncertainty obtained by knowing the weights.
% Distributional uncertainty is mutual information between y and pi at given weights.
% It measures the reduction in label uncertainty obtained by knowing pi, with the weights given.
% In the combined model, this conditional mutual information is also averaged over the weight posterior.
% These expectations use the model distributions at the same input. They do not average over different inputs.
% A deterministic Dirichlet network fixes w at one fitted value, equivalent to a point mass over weights.
% Knowing w then supplies no additional information, so its epistemic term is zero.
% A categorical BNN fixes the class probabilities at f(x;w) once x,w are given.
% In the general hierarchy, this is a point mass for pi at f(x;w), so its distributional term is zero.
% A point mass is a degenerate distribution. The variable can still be written as a random variable mathematically.
% These zeros follow from model assumptions, not evidence that the corresponding source is absent from the problem.
\begin{frame}{Decomposing Predictive Uncertainty}
\label{frm:predictive-entropy}
\label{frm:aleatoric-intuition}
\label{frm:distributional-intuition}
\begin{itemize}
\setlength{\itemsep}{5mm}
\item \textbf{Predictive:}\quad Uncertainty about class label under model's predicted probabilities.
% we use the shannon entropy to measure this. whcih is defined in the appendix.
\item \textbf{Aleatoric:}\quad Label uncertainty that remains once class probabilities are known.
%either f(x;w) for categorical BNN or pi for deterministic Dirichlet network
\item \textbf{Epistemic:}\quad Reduction in label uncertainty obtained by knowing the weights $w$.
%i.e the mutual information between the label and weights
% if we have fixed weights, this implicitly assumes a point mass distribution over the weihgts. I.e. that the weights are known (and are not a random variable with a distribution). Therefore in the dirichlet network, this is zero.
\begin{itemize}
    \item \emph{Deterministic Dirichlet network}: Point mass $P(w=\hat w)=1$.
    \item Assumes no uncertainty about weights. Epistemic uncertainty is zero.\end{itemize}
\item \textbf{Distributional:}\quad Reduction in label uncertainty obtained by knowing $\pi$, given weights $w$.
%i.e the mutual information between the label and the class probabilities given the weights
% if we known class probabilities (f in the case of categorical BNN), this implicitly assumes a point mass distribution over the class probabilities. I.e. that the class probabilities are known (and are not a random variable with a distribution). Therefore in the dirichlet network, this is zero.
\begin{itemize}
    \item \emph{Categorical BNN}: Point mass $P(\pi=f(x;w)|x,w)=1$.
    \item Assumes no uncertainty about the class probabilities. Distributional uncertainty is zero.
\end{itemize}
\item \textbf{Combining weight posterior with a Dirichlet distribution over class probabilities represents all three uncertainty terms.} % and that is the models which we develop in the thesis.
\end{itemize}
\end{frame}




%%%%%%%%%%%%%%%%%%%%%%%%%%%%%%%%%%%%%%%%%%%%%%%%%%%%%%%%%%%%%%%%%%%%%%%%%%%%%%%%

% % Source. Background 2.8 and 2.9, eq:aleatoric-epistemic-decomposition,
% % eq:aleatoric-distributional-split and sec:bg-dirichlet-point-estimate.
% % Speaking notes
% % The categorical BNN represents a posterior distribution over network weights.
% % Given input x and weights w, its class-probability vector is fixed at f(x;w).
% % This is a point mass over class probabilities conditional on x and w, not after averaging over weights.
% % The weights can vary, and each resulting class-probability vector can still allow several labels.
% % The model therefore represents epistemic and aleatoric uncertainty, but its distributional term is zero.
% % A deterministic Dirichlet network instead represents a distribution over class-probability vectors pi.
% % It uses one fitted weight vector, equivalent to a point mass over weights.
% % It therefore represents distributional and aleatoric uncertainty, but its epistemic term is zero.
% % These zeros are structural assumptions. They do not establish that the corresponding source is absent from the problem.
% % We combine the weight posterior with the conditional Dirichlet distribution to represent all three terms.
% \begin{frame}{What Each Model Represents}
% \label{frm:uncertainty-terms}
% \begin{itemize}
% \setlength{\itemsep}{5mm}
% \item \textbf{Categorical BNN}
% \begin{itemize}
% \item Weight posterior $p(w\mid D)$. Knowing $w$ fixes the class probabilities $f(x;w)$ for input $x$.
% \item Distributional term is zero. Represents \textbf{aleatoric and epistemic uncertainty}.
% \end{itemize}
% \item \textbf{Deterministic Dirichlet network}
% \begin{itemize}
% \item $p(\pi\mid x,w)=\mathrm{Dir}(\pi\mid\alpha(x;w))$, but the weights $w$ are fixed.
% \item Epistemic term is zero. Represents \textbf{aleatoric and distributional uncertainty}.
% \end{itemize}
% \item \textbf{Combining both distributions allows all three terms.}
% \end{itemize}
% \end{frame}

% 
